\documentclass[a4paper, 12pt]{article}
\usepackage[brazilian]{babel}
\usepackage{graphicx}
\usepackage{fancyhdr}
\usepackage[left=3cm,right=2cm,top=4cm,bottom=1cm]{geometry}
\usepackage{marvosym}
\usepackage{lipsum}
\usepackage{amsmath}
\usepackage{amssymb}
\usepackage{tcolorbox}
\usepackage{hyperref}
\usepackage{fmtcount}

\newcounter{fis}
\newcounter{mat}

\newcommand{\printfis}{%
  \stepcounter{fis}% incrementa
  FIS\ifnum\value{fis}<10 0\fi% zero à esquerda para <10
  \ifnum\value{fis}<100 0\fi% zero à esquerda para <100
  \arabic{fis}~% imprime número
}

\newcommand{\printmat}{%
  \stepcounter{mat}% incrementa
  MAT\ifnum\value{mat}<10 0\fi% zero à esquerda para <10
  \ifnum\value{mat}<100 0\fi% zero à esquerda para <100
  \arabic{mat}~% imprime número
}

\geometry{headheight = 0.6in, footskip = 1cm, bottom = 2cm}
\fancypagestyle{firstpage}{\fancyhf{}\fancyhead[R]{\includegraphics[height=0.5in, keepaspectratio=true]{Logo_UEL.jpg}}}
\fancypagestyle{plain}{\fancyhf{}\fancyhead[L]{\includegraphics[height=0.5in, keepaspectratio=true]{Logo_UEL.jpg}}\fancyfoot[C]{\thepage}}
\fancypagestyle{empty}{\fancyhf{}\fancyhead[L]{\includegraphics[height=0.5in, keepaspectratio=true]{Logo_UEL.jpg}}\fancyfoot[C]{\thepage}}
\pagestyle{plain}

\title{Núcleo comum do curso de Bacharelado em Física.}
\author{Núcleo Docente Estruturante do curso de Bacharelado em Física.}

\begin{document}

\maketitle

\begin{abstract}
    A Física é uma ciência que tem como objetivo principal a compreensão dos fenômenos
    naturais. Para isso, utilizamos a matemática como linguagem e a experimentação como
    método de validação. Dessa forma, o curso de Bacharelado em Física da Universidade
    Estadual de Londrina tem como objetivo formar profissionais que sejam capazes de
    compreender a linguagem matemática e usá-la para descrever e prever os fenômenos
    naturais. Além disso, o domínio do uso da matemática, aliado a habilidade de mapear
    o mundo real em termos matemáticos, permite que a formação em Física seja uma
    excelente base para a resolução de problemas complexos em diversas áreas do
    conhecimento. Dessa forma, o curso de Bacharelado em Física da Universidade
    Estadual de Londrina é dividido em dois núcleos: o núcleo comum e as enfases. Neste
    documento apresentamos o núcleo comum do curso de Bacharelado em Física da
    Universidade Estadual de Londrina. O núcleo comum é composto por disciplinas
    obrigatórias que todos os alunos devem cursar. O tempo de duração dessa parte do
    curso é de 2 anos.
\end{abstract}

\section*{Núcleo comum}

O objetivo do núcleo comum do curso de Bacharelado em Física da Universidade Estadual
de Londrina é introduzir o aluno ao conhecimento básico da Física e da Matemática.
Nessa parte do curso o aluno terá contato com disciplinas introdutórias de Física e
Matemática onde desenvolverá a fluência matemática combinada com a capacidade de
descrever e modelar elementos da realidade. Tradicionalmente, o curso de física visa
desenvolver essas habilidades com o objetivo de aplicá-las na expansão do conhecimento
científico, seja teórico ou experimental. No entanto, a formação em Física também
permite que o aluno desenvolva habilidades que são úteis em diversas áreas das
atividades humanas, como a capacidade de modelar e resolver problemas complexos, assim
como a habilidade de desenvolver e aplicar técnicas experimentais para validação e
investigação dos modelos propostos.

Por essa razão, o núcleo comum do curso de Bacharelado é composto por disciplinas onde
o aluno terá contato com os conceitos básicos da Física e da Matemática, desenvolvendo
seu entendimento de como essas ideias foram aplicadas de forma bem sucedida. Ou seja,
desse ponto de vista o conteúdo específico das disciplinas serve como um exemplo
concreto de como a matemática pode ser usada para descrever e prever fenômenos
naturais. Evidentemente, para o estudante que deseja seguir carreira acadêmica em
física, o núcleo comum contem os modelos físicos básicos que serão a base para a
compreensão de fenômenos mais complexos. Para o estudante que deseja seguir carreira em
outras áreas, o núcleo comum oferece uma excelente base para a resolução de problemas
complexos em diversas áreas do conhecimento. Para isso, o estudante poderá escolher
entre as diferentes enfases oferecidas pelo curso, que são: Física Industrial, Física
Médica, Bacharelado Tradicional.

O núcleo comum do curso é dividido em dois anos e é composto por disciplinas
semestrais. A seguir apresentamos a grade curricular do núcleo comum do curso e
discutimos os conceitos chave e como eles se relacionam com a ementa das disciplinas e
com o objetivo do curso.

\section{Primeiro ano}

O primeiro ano do curso é um período particularmente desafiador tanto para o corpo
docente quanto para os estudantes. Para os alunos, esse período implica em uma imersão
mais profunda e sistemática nas técnicas matemáticas, além do desafio de conectá-las
aos elementos da realidade que motivaram seu desenvolvimento. Por outro lado, para o
corpo docente, o desafio reside em introduzir conceitos que demandam um esforço
significativo por parte dos estudantes, mantendo-os engajados e capazes de enfrentar os
desafios exigidos.

Para apoiar o processo de aprendizagem, os cursos principais de Cálculo e Física são
acompanhados por disciplinas de Prática, que seguem ementas paralelas focadas na
resolução de exercícios e no treinamento intensivo dos conceitos apresentados. Estas
disciplinas de Prática são ministradas por professores diferentes daqueles responsáveis
pelos cursos teóricos, permitindo uma abordagem complementar e proporcionando aos
alunos diferentes perspectivas sobre os mesmos tópicos. O objetivo dessas práticas é
acompanhar de perto o estudo dos alunos, identificar dificuldades individuais e
coletivas, e proporcionar um ambiente dedicado à consolidação do conhecimento através
da aplicação prática dos conceitos teóricos.

Os dois semestres do primeiro ano são compostos pelas disciplinas abaixo:
\begin{tcolorbox}[title=Primeiro Ano:,colback=blue!5!white,colframe=blue!75!black]
    \begin{tcolorbox}[title=Primeiro semestre:]
        \begin{itemize}
            \item \printmat -- Cálculo I -- Seção~\ref{mat001}
            \item \printmat -- Práticas de Cálculo I -- Seção~\ref{mat002}
            \item \printfis -- Física Geral I -- Seção~\ref{fis001}
            \item \printfis -- Prática de Física Geral I -- Seção~\ref{fis002}
            \item \printfis -- Laboratório de Física Geral I -- Seção~\ref{fis003}
        \end{itemize}
    \end{tcolorbox}
    \begin{tcolorbox}[title=Segundo semestre:]
        \begin{itemize}
            \item \printmat -- Cálculo II -- Seção~\ref{mat003}
            \item \printmat -- Práticas de Cálculo II -- Seção~\ref{mat004}
            \item \printmat -- Métodos Lineares para Física -- Seção~\ref{mat005}
            \item \printfis -- Física Geral IIa -- Seção~\ref{fis004a}
            \item \printfis -- Prática de Física Geral IIa -- Seção~\ref{fis005a}
            \item \printfis -- Física Geral IIb -- Seção~\ref{fis004b}
            \item \printfis -- Prática de Física Geral IIb -- Seção~\ref{fis005b}
            \item \printfis -- Laboratório de Física Geral II -- Seção~\ref{fis006}
        \end{itemize}
    \end{tcolorbox}
\end{tcolorbox}

\subsection*{Grade Horária}

Abaixo apresentamos uma sugestão de grade horária para os dois anos do núcleo comum.
Cada disciplina tem uma carga horária de 60 horas semestrais, exceto as disciplinas
divididas (Física Geral IIa/IIb e Prática de Física Geral IIa/IIb) que têm 30 horas
cada.

\begin{table}[h]
    \centering
    \begin{tabular}{|l|c|c|}
        \hline
        \textbf{Disciplina}            & \textbf{Semestre} & \textbf{Carga Horária} \\
        \hline
        Cálculo I                      & 1º                & 60h                    \\
        Práticas de Cálculo I          & 1º                & 30h                    \\
        Física Geral I                 & 1º                & 60h                    \\
        Prática de Física Geral I      & 1º                & 30h                    \\
        Laboratório de Física Geral I  & 1º                & 60h                    \\
        \hline
        \textbf{Total 1º Semestre}     &                   & \textbf{240h}          \\
        \hline
        \hline
        Cálculo II                     & 2º                & 60h                    \\
        Práticas de Cálculo II         & 2º                & 30h                    \\
        Métodos Lineares para Física   & 2º                & 60h                    \\
        Física Geral IIa               & 2º                & 30h                    \\
        Prática de Física Geral IIa    & 2º                & 30h                    \\
        Física Geral IIb               & 2º                & 30h                    \\
        Prática de Física Geral IIb    & 2º                & 30h                    \\
        Laboratório de Física Geral II & 2º                & 60h                    \\
        \hline
        \textbf{Total 2º Semestre}     &                   & \textbf{300h}          \\
        \hline
        \textbf{Total Primeiro Ano}    &                   & \textbf{660h}          \\
        \hline
    \end{tabular}
    \caption{Grade horária do primeiro ano do núcleo comum.}
\end{table}

Nas próximas seções discutiremos o conteúdo de cada uma dessas disciplinas e como elas
se relacionam com o objetivo do curso. As disciplinas são apresentadas na ordem em que
aparecem no currículo, iniciando com as disciplinas de matemática seguidas pelas de
física em cada semestre.

\pagebreak

\subsection{Cálculo I}
\label{mat001}

O curso de Cálculo I introduz os conceitos fundamentais do cálculo diferencial, que
constituem a base da formulação matemática da Física clássica. Este curso é acompanhado
pela disciplina de Práticas de Cálculo I, onde os conceitos abstratos são aplicados em
contextos geométricos e físicos concretos.

O estudo inicia-se pela construção dos números reais e pela representação de funções de
uma variável real, que permitem descrever relações quantitativas entre grandezas
físicas. A partir dessas relações, são definidos os conceitos de limite e continuidade,
que fornecem o arcabouço necessário para tratar variações contínuas de grandezas
observáveis.

\begin{tcolorbox}[title=Ementa:]
    \begin{itemize}
        \item Conjuntos numéricos.
        \item O conjunto dos números reais.
        \item Intervalos abertos e fechados. Módulo.
        \item Inequações.
        \item Funções de uma variável real.
        \item Composição de funções.
        \item Limite e continuidade de funções reais.
        \item A derivada e derivação.
        \item Propriedades (linearidade).
        \item Regra da cadeia.
        \item Derivação implícita.
        \item Regra de L'Hôpital.
        \item Fórmula e polinômio de Taylor.
        \item Funções reais de várias variáveis: derivadas parciais.
        \item Derivadas parciais e regra da cadeia.
        \item Máximos e mínimos.
    \end{itemize}
\end{tcolorbox}

\subsection{Práticas de Cálculo I}
\label{mat002}

A disciplina de Práticas de Cálculo I funciona em paralelo ao curso de Cálculo I,
proporcionando aos alunos um espaço dedicado para consolidar os conceitos abstratos do
cálculo através de sua aplicação em contextos geométricos e físicos concretos. Esta
disciplina é fundamental para estabelecer a ponte entre a matemática formal apresentada
no curso teórico e suas aplicações práticas.

O curso começa com a representação geométrica, explorando retas e coordenadas no plano
cartesiano. Os alunos desenvolvem a habilidade de visualizar funções e relações
matemáticas geometricamente, uma competência essencial para a compreensão intuitiva dos
conceitos do cálculo. A introdução às cônicas e quádricas expande essa compreensão para
curvas e superfícies mais complexas, preparando o terreno para estudos futuros em
geometria diferencial.

O estudo de diferentes sistemas de coordenadas é particularmente importante, pois muitos
problemas físicos são mais facilmente resolvidos quando descritos no sistema de
coordenadas apropriado. Os alunos aprendem a trabalhar com coordenadas cartesianas,
polares, cilíndricas e esféricas, desenvolvendo a flexibilidade necessária para
escolher a representação mais adequada para cada problema.

A aplicação do cálculo ao movimento acelerado unidimensional proporciona aos alunos uma
primeira experiência concreta de como as derivadas descrevem mudanças instantâneas em
sistemas físicos. Problemas unidimensionais servem como introdução à modelagem
matemática de fenômenos físicos, onde os conceitos de velocidade e aceleração ganham
significado tangível através das derivadas.

O conceito de vetores covariantes e contravariantes na geometria euclidiana introduz os
alunos a uma perspectiva mais sofisticada sobre objetos geométricos, preparando-os para
estudos futuros em física teórica e geometria diferencial. Finalmente, o estudo de
campos escalares, como distribuição de calor e densidade, conecta o cálculo a fenômenos
físicos reais, mostrando como funções de várias variáveis podem descrever propriedades
que variam no espaço.

\begin{tcolorbox}[title=Ementa:]
    \begin{itemize}
        \item Retas e coordenadas. Plano cartesiano.
        \item Introdução às cônicas e quádricas.
        \item Sistemas de coordenadas.
        \item Movimento acelerado unidimensional.
        \item Problemas unidimensionais com derivadas.
        \item Vetores covariantes e contravariantes na geometria euclidiana.
        \item Geometria euclidiana.
        \item Campos escalares (distribuição de calor, densidade, etc.).
    \end{itemize}
\end{tcolorbox}

\subsection{Física Geral I}
\label{fis001}

O curso de Física Geral I introduz três conceitos-chave da Física clássica: o modelo de
partícula, sua cinemática e dinâmica. Esse curso é acompanhado pela disciplina de
Prática de Física Geral I, onde os alunos consolidam o aprendizado através da resolução
de problemas e exercícios, com foco na aplicação dos conceitos teóricos em situações
concretas. O modelo de partícula pontual é de suma
importância por diversas razões. Primeiramente, fornece um modelo simples que descreve
o movimento de objetos macroscópicos. Em segundo lugar, serve como base para a
descrição de sistemas mais complexos, como corpos rígidos. Além de ser um modelo
aproximado, é a peça fundamental para a descrição de sistemas mais complexos, os quais
passam a ser representados como um conjunto de partículas.

A cinemática desse modelo envolve a descrição do movimento de partículas, que é
definido pelo espaço e pelo tempo. Essa descrição ocorre de duas maneiras: uma física,
onde determinamos o que podemos medir, e uma matemática, onde definimos os objetos
matemáticos que descrevem essas medidas. Em outras palavras, a partir da observação dos
pontos materiais, estabelecemos um sistema de coordenadas e associamos a cada ponto as
medidas físicas de distância a partir de um ponto de referência. Item da ementa: {\bf
Medidas e grandezas em Física.}

De forma prática, o ambiente é descrito em um espaço tridimensional, onde cada ponto é
associado a um conjunto de coordenadas representando medidas físicas de distância.
Assim, a cinemática demanda a definição da geometria desse espaço. Uma vez definido,
podemos descrever o movimento de uma partícula em termos de sua posição e velocidade.
Em termos matemáticos, as posições são representadas como vetores em um espaço
tridimensional $\mathbb{M} = \mathbb{R}^3$. Além disso, o estado físico de uma
partícula também requer a descrição de sua velocidade, que, novamente, pode ser
modelada matematicamente como um vetor em outro espaço tridimensional $\mathbb{V} =
    \mathbb{R}^3$. Portanto, compreender a cinemática exige habilidade no trabalho com
vetores e sistemas de coordenadas. Item da ementa: {\bf Vetores e sistemas de
coordenadas}.

Assim, a cinemática trata da descrição do movimento de partículas em um espaço
tridimensional, independentemente das causas desse movimento. Em essência, fornece um
modelo matemático para descrever trajetórias, velocidades e acelerações de partículas
na ausência explícita de forças. Inclui o estudo de movimentos retilíneos uniformes e
uniformemente acelerados, bem como de movimentos circulares, que podem ser
interpretados como casos particulares observados em referenciais acelerados. Sob uma
perspectiva mais ampla, o movimento livre, isto é, na ausência de interações, pode ser
entendido como o movimento ao longo de uma geodésica\footnote{Em termos gerais, uma
    geodésica é o caminho de menor ``distância'' (ou tempo próprio) em um espaço ou
    espaço-tempo, conceito fundamental em teorias como a relatividade geral.}.  Item da
ementa: {\bf Cinemática da partícula}.

O próximo elemento a ser introduzido é a dinâmica, que descreve o movimento de
partículas na presença de forças. No entanto, ao contrário da cinemática, a dinâmica
não possui um caráter descritivo universal. Isso significa que, dependendo das
características específicas do sistema, diferentes leis podem ser aplicadas. Portanto,
por meio das leis de Newton, introduzimos o conceito de força para modelar como as
partículas respondem a essas interações.

Em termos matemáticos, a dinâmica é descrita por meio de equações diferenciais que
delineiam como as forças influenciam o movimento das partículas e como estas reagem a
tais interações. É importante notar que a dinâmica não apenas introduz o conceito de
força, mas também explora como essas forças se relacionam com a massa e a aceleração
das partículas, conforme expresso pelas famosas leis de Newton. Essas leis constituem
os princípios fundamentais da dinâmica e servem como alicerce para a modelagem de uma
vasta gama de fenômenos físicos, desde o movimento dos planetas até o funcionamento de
máquinas e sistemas mecânicos. A dinâmica pode ser dividida em duas partes: os
princípios que governam as interações entre partículas (ação e reação, forças centrais,
etc.) e o estudo da evolução temporal desses sistemas. Item da ementa: {\bf Os
princípios da dinâmica}.

A construção do sistema de referência descrita acima depende da escolha de um ponto de
referência a partir do qual as medidas são feitas. Dependendo do movimento do sistema,
esses referenciais podem ser classificados como inerciais ou não inerciais. O conceito
de referencial inercial é fundamental para a descrição do movimento de partículas, uma
vez que fornece um quadro de referência no qual temos regras bem definidas para
descrever o movimento de partículas. Em contraste, os referenciais não inerciais são
aqueles que estão em movimento acelerado em relação a um referencial inercial. Esses
referenciais podem ser usados para descrever o movimento de partículas, mas exigem a
introdução de forças fictícias para manter a consistência com as leis de Newton. Item
da ementa: {\bf Referenciais inerciais e não-inerciais}.

Quando exploramos a dinâmica de um sistema, podemos categorizar as forças em dois tipos
distintos: conservativas e dissipativas. Nas forças conservativas, destacamos a
capacidade de definir uma grandeza chamada energia, a qual é preservada ao longo do
tempo. Além disso, a estrutura das forças conservativas nos permite estabelecer outra
grandeza fundamental, denominada momento linear, que também é mantida inalterada ao
longo da evolução temporal do sistema. Em resumo, os princípios físicos e geométricos
subjacentes às forças possibilitam a definição de grandezas conservadas que desempenham
um papel crucial na descrição e compreensão do comportamento dinâmico dos sistemas.
Item da ementa: {\bf Conservação da energia e momento linear}.

Uma progressão natural do conceito de conservação reside na interconexão entre
simetrias e leis de conservação. De maneira abrangente, as simetrias presentes em um
sistema físico estão intrinsecamente ligadas às leis de conservação que governam esse
sistema. Em outras palavras, a existência de determinadas simetrias em um sistema
implica na existência de leis de conservação que regem suas dinâmicas. Assim, a relação
entre simetrias e leis de conservação é um princípio fundamental na compreensão de como
as leis físicas são derivadas a partir de princípios mais fundamentais, proporcionando
uma visão mais profunda sobre a estrutura e o funcionamento dos fenômenos físicos. Item
da ementa: {\bf Simetrias e leis de conservação}.

Acima, delineamos os conceitos fundamentais abordados no curso de Física Geral I. Para
ilustrar esses conceitos, o curso apresenta exemplos concretos de sua aplicação em
diversos contextos. Um desses exemplos destacados é a gravitação, um fenômeno
amplamente compreendido e descrito pela Física Clássica. Item da ementa: {\bf
Gravitação}.

Além disso, o modelo do ponto material é a base para a descrição de sistemas
mais complexos, como corpos rígidos. A cinemática e a dinâmica dos corpos rígidos
são introduzidas no curso de Física Geral I, onde o aluno aprenderá a descrever o
movimento de corpos rígidos em termos de suas posições e velocidades.

\begin{tcolorbox}[title=Ementa:]
    \begin{itemize}
        \item Vetores e sistemas de coordenadas.
        \item Cinemática da partícula.
        \item Os princípios da dinâmica.
        \item Referenciais inerciais e não-inerciais.
        \item Conservação da energia e momento linear.
        \item Simetrias e leis de conservação.
    \end{itemize}
\end{tcolorbox}

\subsection{Prática de Física Geral I}
\label{fis002}

A disciplina de Prática de Física Geral I é um componente essencial do primeiro
semestre, funcionando em paralelo ao curso de Física Geral I. O objetivo desta
disciplina é proporcionar aos alunos a oportunidade de consolidar os conceitos teóricos
através da resolução de problemas práticos e exercícios, desenvolvendo habilidades
fundamentais para o estudo da Física.

Um aspecto crucial desta disciplina é a construção de sistemas de coordenadas a partir
de medidas físicas reais. Os alunos aprendem a estabelecer a conexão entre o mundo
físico observável e sua representação matemática, compreendendo como transformar
observações experimentais em descrições quantitativas. Esta habilidade é fundamental
para toda a formação em Física e representa o primeiro passo na modelagem de fenômenos
naturais.

A disciplina também enfatiza a conversão entre diferentes grandezas e a familiarização
com o Sistema Internacional de Unidades (SI). Os alunos desenvolvem a capacidade de
trabalhar com diferentes escalas e sistemas de medida, uma habilidade essencial para a
comunicação científica e para a resolução de problemas reais.

A análise dimensional é outro tópico central desta disciplina. Esta técnica poderosa
permite verificar a consistência de equações, estimar ordens de grandeza e compreender
as relações entre diferentes quantidades físicas. É uma ferramenta que os alunos
utilizarão durante toda sua formação e carreira profissional.

Os problemas envolvendo referenciais com aceleração constante permitem aos alunos
aplicar os conceitos de cinemática em situações concretas, desenvolvendo intuição sobre
movimento relativo e transformações entre sistemas de referência. Finalmente, a
disciplina explora problemas envolvendo simetrias e quantidades conservadas,
estabelecendo a conexão profunda entre princípios de simetria e leis de conservação, um
dos conceitos mais fundamentais e elegantes da Física.

\begin{tcolorbox}[title=Ementa:]
    \begin{itemize}
        \item Construção de sistemas de coordenadas a partir de medidas físicas.
        \item Conversão entre grandezas, sistema SI.
        \item Análise dimensional.
        \item Problemas com referenciais com aceleração constante.
        \item Problemas com simetrias e quantidades conservadas.
    \end{itemize}
\end{tcolorbox}

\subsection{Laboratório de Física Geral I}
\label{fis003}

A disciplina de Laboratório de Física Geral I complementa os cursos teóricos e de
prática, proporcionando aos alunos experiência direta com experimentos físicos que
ilustram os conceitos de mecânica estudados em Física Geral I.

\begin{tcolorbox}[title=Ementa:]
    \begin{itemize}
        \item A ementa será definida para complementar os tópicos de mecânica.
    \end{itemize}
\end{tcolorbox}

\subsection{Cálculo II}
\label{mat003}

O curso de Cálculo II introduz o cálculo integral, complementando os conceitos de
cálculo diferencial apresentados em Cálculo I. Este curso é acompanhado pela disciplina
de Práticas de Cálculo II, onde os alunos aplicam as técnicas de integração a problemas
físicos e geométricos concretos, além de serem introduzidos às equações diferenciais
separáveis.

A integral de Riemann é apresentada inicialmente através de sua definição geométrica,
permitindo aos alunos desenvolver intuição sobre o conceito de acumulação e área sob
curvas. A definição analítica da integral consolida essa compreensão, estabelecendo as
bases rigorosas para o cálculo integral.

O estudo das funções transcendentes e suas integrais expande o repertório de funções
que os alunos podem manipular, incluindo funções exponenciais, logarítmicas e
trigonométricas. As técnicas de integração são apresentadas como ferramentas essenciais
para resolver integrais mais complexas, desenvolvendo a habilidade dos alunos de
reconhecer padrões e escolher métodos apropriados.

As integrais impróprias estendem o conceito de integração para domínios ilimitados ou
funções com descontinuidades, enquanto as integrais múltiplas generalizam a integração
para funções de várias variáveis, preparando os alunos para aplicações em física que
envolvem cálculos em duas ou três dimensões.

\begin{tcolorbox}[title=Ementa:]
    \begin{itemize}
        \item A integral de Riemann.
        \item Definição analítica da integral.
        \item Estudo das funções transcendentes.
        \item Fórmulas de integrais conhecidas.
        \item Técnicas de integração.
        \item Integrais impróprias.
        \item Integrais múltiplas.
    \end{itemize}
\end{tcolorbox}

\subsection{Práticas de Cálculo II}
\label{mat004}

A disciplina de Práticas de Cálculo II acompanha o curso de Cálculo II, focando na
aplicação das técnicas de integração a problemas concretos e na introdução às equações
diferenciais. Esta disciplina é essencial para que os alunos desenvolvam fluência no
uso das integrais e compreendam seu significado físico e geométrico.

O curso inicia com aplicações imediatas de integrais em uma dimensão, onde os alunos
exploram o cálculo de áreas, volumes, trabalho, energia e outras grandezas físicas que
podem ser expressas como integrais. Estas aplicações demonstram o poder do cálculo
integral para resolver problemas práticos e estabelecem conexões diretas com fenômenos
físicos observáveis.

Um componente crucial desta disciplina é a introdução às equações diferenciais
separáveis, particularmente aquelas do tipo $dy/dx = f(x)g(y)$. Esta classe de equações
aparece naturalmente em inúmeros contextos físicos, desde crescimento populacional até
decaimento radioativo e circuitos elétricos. Os alunos aprendem a reconhecer, resolver
e interpretar estas equações, desenvolvendo uma compreensão de como as taxas de
variação determinam o comportamento dinâmico de sistemas.

As aplicações imediatas de integrais em várias dimensões expandem o repertório dos
alunos para problemas tridimensionais e além. Conceitos como densidade, centro de
massa, fluxo através de superfícies e outros fenômenos que requerem integração múltipla
são explorados, preparando os alunos para cursos mais avançados de física e matemática
aplicada.

Esta disciplina desempenha um papel fundamental na formação do físico, pois é através
da resolução de problemas concretos que os alunos desenvolvem a intuição necessária
para aplicar o cálculo integral em contextos diversos, além de adquirir as primeiras
ferramentas para modelar sistemas dinâmicos através de equações diferenciais.

\begin{tcolorbox}[title=Ementa:]
    \begin{itemize}
        \item Aplicações imediatas de integrais em uma dimensão.
        \item Algumas equações diferenciais separáveis.
        \item Equações do tipo $dy/dx=f(x)g(y)$.
        \item Aplicações imediatas de integrais em várias dimensões.
    \end{itemize}
\end{tcolorbox}

\subsection{Métodos Lineares para Física}
\label{mat005}

A disciplina de Métodos Lineares para Física introduz os conceitos fundamentais de
álgebra linear com foco em suas aplicações na física. Este curso é essencial para o
desenvolvimento das habilidades matemáticas necessárias para compreender sistemas
físicos descritos por equações lineares.

Os alunos aprendem sobre vetores e espaços vetoriais, combinacões lineares e bases,
conceitos fundamentais para a descrição de sistemas físicos. O estudo de transformações
lineares e suas representações matriciais fornece ferramentas para analisar como
operadores físicos atuam sobre estados do sistema.

Autovetores e autovalores são apresentados como conceitos centrais para a análise de
sistemas físicos, incluindo modos normais de vibração e estados estacionários em
mecânica quântica. A diagonalização de matrizes simétricas e operadores ortogonais são
estudados com aplicações a rotações e outros problemas físicos.

\begin{tcolorbox}[title=Ementa:]
    \begin{itemize}
        \item Vetores e espaços vetoriais usados em física.
        \item Combinações lineares, subespaços e bases.
        \item Produto interno, norma e projeções.
        \item Transformações lineares e representações matriciais.
        \item Operações com matrizes: soma, produto, inversa.
        \item Propriedades principais: determinante, traço, posto.
        \item Mudança de base e interpretações físicas simples.
        \item Autovalores e autovetores.
        \item Diagonalização no caso simétrico.
        \item Operadores ortogonais e aplicações a rotações.
        \item Modos próprios em sistemas lineares simples.
    \end{itemize}
\end{tcolorbox}

\subsection{Física Geral II}
\label{fis004a}
\label{fis004b}

O curso de Física Geral II expande o conhecimento adquirido em Física Geral I,
introduzindo dois novos domínios fundamentais da Física clássica: oscilações e ondas, e
termodinâmica. Este curso é dividido em duas partes complementares que exploram
fenômenos distintos mas igualmente fundamentais para a compreensão da natureza.

Na primeira parte do curso (Física IIa), os alunos são introduzidos ao estudo de
oscilações e ondas. O estudo começa com oscilações livres e forçadas, fenômenos
ubíquos na natureza que descrevem desde o movimento de pêndulos até vibrações
moleculares. O conceito de ressonância é particularmente importante, pois aparece em
inúmeros contextos físicos e tem aplicações práticas significativas em engenharia e
tecnologia.

A propagação de ondas e o conceito de velocidade de fase são introduzidos como
extensões naturais do estudo de oscilações. Os alunos aprendem como perturbações se
propagam através de diferentes meios e como descrever matematicamente esse fenômeno. Os
conceitos de índice de refração, difração, reflexão e transmissão são explorados no
contexto da propagação de ondas, culminando com o princípio de Huygens e uma introdução
à óptica geométrica.

Na segunda parte do curso (Física IIb), o foco se volta para a física de fluidos e
termodinâmica. Os alunos estudam as propriedades dos fluidos, conceitos de pressão e o
princípio de Arquimedes. A termodinâmica é introduzida através dos conceitos de
temperatura e calor, seguidos pela primeira e segunda leis da termodinâmica, que
estabelecem os princípios fundamentais de conservação e transformação de energia.

O curso também apresenta a teoria cinética dos gases, que estabelece a ponte entre a
descrição microscópica (molecular) e macroscópica (termodinâmica) da matéria. A
distribuição de Maxwell-Boltzmann é introduzida como uma ferramenta fundamental para
compreender o comportamento estatístico de sistemas com muitas partículas.

\begin{tcolorbox}[title=Ementa Física IIa:]
    \begin{itemize}
        \item Oscilações livres e forçadas.
        \item Ressonância.
        \item Propagação e velocidade de fase.
        \item Índice de refração e difração.
        \item Reflexão e transmissão.
        \item Difração e o princípio de Huygens.
        \item Óptica geométrica.
    \end{itemize}
\end{tcolorbox}

\begin{tcolorbox}[title=Ementa Física IIb:]
    \begin{itemize}
        \item Propriedades dos fluidos.
        \item Pressão num fluido.
        \item Princípio de Arquimedes.
        \item Temperatura e calor.
        \item Primeira e segunda leis da termodinâmica.
        \item Teoria cinética dos gases.
        \item Distribuição de Maxwell-Boltzmann.
    \end{itemize}
\end{tcolorbox}

\subsection{Prática de Física Geral II}
\label{fis005a}
\label{fis005b}

A disciplina de Prática de Física Geral II acompanha o curso de Física Geral II e tem
como objetivo proporcionar aos alunos experiência prática com os conceitos de
oscilações, ondas e termodinâmica. Similar à estrutura do curso teórico, as práticas
são divididas em duas partes correspondentes.

As Práticas de Física IIa focam no estudo experimental de sistemas oscilatórios. O
pêndulo simples serve como um sistema modelo fundamental, permitindo aos alunos
observar e quantificar movimento harmônico simples, medir períodos de oscilação e
verificar as leis que governam esse tipo de movimento. O estudo se estende ao pêndulo
amortecido, onde os alunos investigam como a dissipação de energia afeta o
comportamento oscilatório, observando fenômenos como o decaimento exponencial da
amplitude e as mudanças na frequência de oscilação.

Essas práticas permitem que os alunos desenvolvam habilidades experimentais importantes,
como a coleta e análise de dados, o ajuste de modelos teóricos a dados experimentais, e
a compreensão das limitações dos modelos idealizados quando confrontados com situações
reais onde atrito e outras forças dissipativas estão presentes.

\begin{tcolorbox}[title=Ementa Prática de Física IIa:]
    \begin{itemize}
        \item Pêndulo simples.
        \item Pêndulo amortecido.
    \end{itemize}
\end{tcolorbox}

\begin{tcolorbox}[title=Ementa Prática de Física IIb:]
    \begin{itemize}
        \item A ementa das práticas de Física IIb será definida para complementar os
              tópicos de fluidos e termodinâmica.
    \end{itemize}
\end{tcolorbox}

\subsection{Laboratório de Física Geral II}
\label{fis006}

A disciplina de Laboratório de Física Geral II complementa as práticas, fornecendo uma
experiência mais aprofundada com experimentos físicos relacionados aos tópicos do curso
teórico.

\pagebreak

\section{Segundo Ano}
\label{segundoano}

No segundo ano do curso, o aluno terá contato com disciplinas abaixo:
\begin{tcolorbox}[title=Segundo Ano:,colback=blue!5!white,colframe=blue!75!black]
    \begin{tcolorbox}[title=Terceiro semestre:]
        \begin{itemize}
            \item \printmat -- Cálculo III -- Seção~\ref{mat006}
            \item \printmat -- Práticas de Cálculo III -- Seção~\ref{mat007}
            \item \printfis -- Computação para Física -- Seção~\ref{fis007}
            \item \printfis -- Física Geral III -- Seção~\ref{fis008}
            \item \printfis -- Prática de Física Geral III -- Seção~\ref{fis009}
            \item \printfis -- Laboratório de Física Geral III -- Seção~\ref{fis010}
        \end{itemize}
    \end{tcolorbox}
    \begin{tcolorbox}[title=Quarto semestre:]
        \begin{itemize}
            \item \printmat -- Cálculo IV -- Seção~\ref{mat008}
            \item \printmat -- Práticas de Cálculo IV -- Seção~\ref{mat009}
            \item \printmat -- Dinâmica de Sistemas Físicos -- Seção~\ref{mat010}
            \item \printfis -- Física Geral IV -- Seção~\ref{fis011}
            \item \printfis -- Prática de Física Geral IV -- Seção~\ref{fis012}
            \item \printfis -- Laboratório de Física Geral IV -- Seção~\ref{fis013}
        \end{itemize}
    \end{tcolorbox}
\end{tcolorbox}

\subsection*{Grade Horária}

\begin{table}[h]
    \centering
    \begin{tabular}{|l|c|c|}
        \hline
        \textbf{Disciplina}             & \textbf{Semestre} & \textbf{Carga Horária} \\
        \hline
        Cálculo III                     & 3º                & 60h                    \\
        Práticas de Cálculo III         & 3º                & 30h                    \\
        Computação para Física          & 3º                & 60h                    \\
        Física Geral III                & 3º                & 60h                    \\
        Prática de Física Geral III     & 3º                & 30h                    \\
        Laboratório de Física Geral III & 3º                & 60h                    \\
        \hline
        \textbf{Total 3º Semestre}      &                   & \textbf{300h}          \\
        \hline
        \hline
        Cálculo IV                      & 4º                & 60h                    \\
        Práticas de Cálculo IV          & 4º                & 30h                    \\
        Dinâmica de Sistemas Físicos    & 4º                & 60h                    \\
        Física Geral IV                 & 4º                & 60h                    \\
        Prática de Física Geral IV      & 4º                & 30h                    \\
        Laboratório de Física Geral IV  & 4º                & 60h                    \\
        \hline
        \textbf{Total 4º Semestre}      &                   & \textbf{300h}          \\
        \hline
        \textbf{Total Segundo Ano}      &                   & \textbf{600h}          \\
        \hline
        \textbf{Total Núcleo Comum}     &                   & \textbf{1380h}         \\
        \hline
    \end{tabular}
    \caption{Grade horária do segundo ano do núcleo comum.}
\end{table}

\subsection{Cálculo III}
\label{mat006}

O curso de Cálculo III estende os conceitos de cálculo para funções de várias
variáveis e introduz o cálculo vetorial.

\begin{tcolorbox}[title=Ementa:]
    \begin{itemize}
        \item A ementa será definida.
    \end{itemize}
\end{tcolorbox}

\subsection{Práticas de Cálculo III}
\label{mat007}

A disciplina de Práticas de Cálculo III acompanha o curso de Cálculo III, focando na
aplicação prática dos conceitos de cálculo vetorial e funções de várias variáveis.

\begin{tcolorbox}[title=Ementa:]
    \begin{itemize}
        \item A ementa será definida.
    \end{itemize}
\end{tcolorbox}

\subsection{Computação para Física}
\label{fis007}

Esta disciplina introduz os alunos às técnicas computacionais aplicadas à resolução de
problemas físicos. O curso foca no desenvolvimento de habilidades práticas em
programção científica, utilizando Python como linguagem principal.

Os alunos aprendem a organizar projetos científicos, manipular e analisar dados
experimentais, e utilizar bibliotecas fundamentais como NumPy, SciPy e Matplotlib. O
curso aborda métodos numéricos básicos para resolução de problemas físicos, incluindo
zeros de funções, integração numérica e solução de equações diferenciais ordinárias.

\begin{tcolorbox}[title=Ementa:]
    \begin{itemize}
        \item Organização básica de arquivos, ambientes e fluxo de trabalho para projetos em física.
        \item Introdução à programação científica em Python: variáveis, funções, estruturas simples.
        \item Uso de notebooks interativos para cálculos e visualização.
        \item Bibliotecas fundamentais: NumPy, SciPy, Matplotlib.
        \item Manipulação e análise de dados experimentais simples.
        \item Resolução numérica de problemas elementares: zeros de funções, integrais, EDOs simples.
        \item Representação e ajustes de dados experimentais: mínimos quadrados básicos.
        \item Automatização de tarefas e reprodutibilidade.
        \item Boas práticas de programação para projetos científicos.
    \end{itemize}
\end{tcolorbox}

\subsection{Física Geral III}
\label{fis008}

O curso de Física Geral III é um passo importante na exploração dos princípios
fundamentais da Física clássica, concentrando-se no eletromagnetismo, uma das quatro
forças fundamentais da natureza. Seu objetivo primordial é duplo: compreender tanto a
base empírica das leis eletromagnéticas quanto sua formulação matemática. Inicialmente,
os alunos exploram as leis do eletromagnetismo de forma fenomenológica, ou seja,
compreendendo seus efeitos observáveis na natureza. Posteriormente, avançam para a
unificação dessas leis em um único formalismo, exemplificado pelas equações de Maxwell.

Nesse curso, os alunos são apresentados ao eletromagnetismo de forma simplificada,
utilizando modelos com soluções analíticas e simétricas. O objetivo é que compreendam a
natureza das interações eletromagnéticas e como elas se manifestam em diferentes
contextos, sem a introdução de formalismos matemáticos complexos nesta etapa inicial.
Os formalismos mais avançados são reservados para cursos posteriores, como o de
Eletromagnetismo I, onde os alunos podem avançar para uma formulação mais rigorosa e
abrangente das leis do eletromagnetismo.

O curso é dividido em duas partes complementares: Física Geral IIIa e Física Geral
IIIb. Na primeira parte, os conceitos fundamentais do eletromagnetismo são
introduzidos de forma fenomenológica. Na segunda parte, o formalismo matemático é
aprofundado através da revisão de cálculo vetorial e da formulação das equações de
Maxwell.

No curso de Física Geral I (ver seção~\ref{fis001}), assim como no presente curso, o
modelo de partícula é fundamental para descrever o eletromagnetismo. Contudo, o
eletromagnetismo introduz um novo conceito crucial: a carga elétrica. Essa propriedade
intrínseca da matéria é representada matematicamente como um escalar, enriquecendo o
modelo de partícula com essa característica fundamental. Sua introdução ocorre de forma
fenomenológica, sendo apresentada como a força que governa a interação entre partículas
carregadas, conforme descrito pela lei de Coulomb, derivada de observações
experimentais. A partir das características dessa interação, exploramos o conceito de
campo elétrico e suas propriedades, incluindo a Lei de Gauss e o potencial elétrico.
Item da ementa: {\bf Lei de Coulomb e campo elétrico; Lei de Gauss; Potencial elétrico}.

\begin{tcolorbox}[title=Ementa Física Geral IIIa:]
    \begin{itemize}
        \item Lei de Coulomb e campo elétrico.
        \item Lei de Gauss.
        \item Potencial elétrico.
        \item Capacitores e dielétricos.
        \item Corrente, resistência e força eletromotriz.
        \item Circuitos e instrumentos de corrente contínua.
        \item Campo magnético de uma corrente.
        \item Forças magnéticas sobre portadores de correntes.
        \item Força eletromotriz induzida.
        \item Circuitos de corrente alternada.
        \item Propriedades magnéticas da matéria.
        \item Correntes alternadas.
    \end{itemize}
\end{tcolorbox}

\begin{tcolorbox}[title=Ementa Física Geral IIIb:]
    \begin{itemize}
        \item Revisão de cálculo vetorial: álgebra vetorial, produtos.
        \item Cálculo diferencial, os operadores grad, div e rot, derivadas segundas.
        \item Cálculo integral e teoremas integrais.
        \item Equações de Maxwell.
        \item Oscilações eletromagnéticas.
        \item Ondas eletromagnéticas.
        \item Polarização.
        \item Interferência e difração da luz.
    \end{itemize}
\end{tcolorbox}

\subsection{Prática de Física Geral III}
\label{fis009}

A disciplina de Prática de Física Geral III acompanha o curso de Física Geral III,
proporcionando aos alunos a oportunidade de consolidar os conceitos de eletromagnetismo
através da resolução de problemas práticos e exercícios.

\begin{tcolorbox}[title=Ementa:]
    \begin{itemize}
        \item A ementa será definida para complementar os tópicos de eletromagnetismo.
    \end{itemize}
\end{tcolorbox}

\subsection{Laboratório de Física Geral III}
\label{fis010}

A disciplina de Laboratório de Física Geral III complementa o curso teórico,
proporcionando experiência prática com experimentos de eletromagnetismo.

\begin{tcolorbox}[title=Ementa:]
    \begin{itemize}
        \item A ementa será definida para complementar os tópicos de eletromagnetismo.
    \end{itemize}
\end{tcolorbox}

\subsection{Cálculo IV}
\label{mat008}

O curso de Cálculo IV aprofunda o estudo de cálculo vetorial, incluindo integrais de
linha e superfície, e os teoremas fundamentais do cálculo vetorial.

\begin{tcolorbox}[title=Ementa:]
    \begin{itemize}
        \item A ementa será definida.
    \end{itemize}
\end{tcolorbox}

\subsection{Práticas de Cálculo IV}
\label{mat009}

A disciplina de Práticas de Cálculo IV acompanha o curso de Cálculo IV, focando na
aplicação prática dos conceitos de integrais de linha e superfície e dos teoremas
fundamentais do cálculo vetorial.

\begin{tcolorbox}[title=Ementa:]
    \begin{itemize}
        \item A ementa será definida.
    \end{itemize}
\end{tcolorbox}

\subsection{Dinâmica de Sistemas Físicos}
\label{mat010}

Esta disciplina introduz a teoria e métodos de solução de equações diferenciais
ordinárias aplicadas a sistemas físicos. O curso enfatiza a compreensão do comportamento
dinâmico de sistemas através de modelos matemáticos.

Os alunos estudam modelos de primeira ordem (decaimento, crescimento, relaxação) e o
oscilador harmônico em suas várias formas: livre, amortecido e forçado. O conceito de
ressonância é explorado em profundidade. Sistemas acoplados e modos normais são
introduzidos, conectando-se com os conceitos de álgebra linear estudados anteriormente.

\begin{tcolorbox}[title=Ementa:]
    \begin{itemize}
        \item Leis de evolução e taxas de variação em sistemas físicos simples.
        \item Modelos de 1ª ordem: decaimento, crescimento, relaxação.
        \item Equações de movimento para forças lineares e atrito viscoso.
        \item Oscilador harmônico: solução, energia e representação em fase.
        \item Oscilador amortecido: regimes de amortecimento.
        \item Oscilador forçado: resposta transitória, regime estacionário e ressonância.
        \item Sistemas acoplados de poucos graus de liberdade; modos normais.
        \item Representação matricial da dinâmica linear; autovalores e frequências naturais.
        \item Estabilidade e comportamento de equilíbrios em modelos físicos simples.
        \item Resposta a forçamentos periódicos; decomposição harmônica.
    \end{itemize}
\end{tcolorbox}

\subsection{Física Geral IV}
\label{fis011}

O curso de Física Geral IV aborda tópicos avançados da física clássica e introduz
conceitos da física moderna.

\begin{tcolorbox}[title=Ementa:]
    \begin{itemize}
        \item A ementa será definida.
    \end{itemize}
\end{tcolorbox}

\subsection{Prática de Física Geral IV}
\label{fis012}

A disciplina de Prática de Física Geral IV acompanha o curso de Física Geral IV,
proporcionando aos alunos a oportunidade de consolidar os conceitos avançados de física
através da resolução de problemas práticos e exercícios.

\begin{tcolorbox}[title=Ementa:]
    \begin{itemize}
        \item A ementa será definida.
    \end{itemize}
\end{tcolorbox}

\subsection{Laboratório de Física Geral IV}
\label{fis013}

A disciplina de Laboratório de Física Geral IV complementa o curso teórico com
experimentos relacionados aos tópicos avançados de física.

\begin{tcolorbox}[title=Ementa:]
    \begin{itemize}
        \item A ementa será definida.
    \end{itemize}
\end{tcolorbox}


\end{document}
